 \section{Proof of Theorem~\ref{thm:GHpert}}\label{sec:proof}

 The proof of Theorem~\ref{thm:GHpert} starts by establishing an inversion result that will later be applied to $I+A(\lambda)E^\pm$, where $E^\pm$ are the Green operators of $P$. This is proved first in Sobolev spaces $H^s_K$, using the analytic Fredholm theorem, and then extended to $C^\infty_K$. The main part of the proof then uses this information to define Green operators for $P+ A(\lambda)$.


 \begin{Theorem}\label{thm:Sobolevinversion}
 	Let $K\subset M$ be a fixed topologically regular compact set. Suppose
 	$Y(\lambda)$, $\lambda\in\CC$, is a~family of linear maps
 	$Y(\lambda)\colon\DD_K'\to\DD'_K$, each restricting to continuous maps $Y_s(\lambda)\colon H^s_K\to H^{s+\delta}_K$ for all $s\ge s_*\in\RR$ and some fixed $\delta>0$. Assume also that $Y(0)=0$.
 	Let $\hat{Y}_s(\lambda)\in\LL(H^s_K)$ denote the compact maps obtained by composing $Y_s(\lambda)$ with the embedding $H^{s+\delta}_K\hookrightarrow H^s_K$. Suppose that $\lambda \to \hat{Y}_s(\lambda)$ is holomorphic on $\CC$ in $\LL_b(H^s_K)$ for all $s\ge s_*$. Then
 	\begin{enumerate}\itemsep=0pt
 		\item[$(a)$] The set
 		\begin{equation}\label{eq:Sdef}
 			S=\big\{\lambda\in\CC\colon \ker (I+ Y(\lambda))|_{C^\infty_K}\neq 0\big\}
 		\end{equation}
 		is a discrete subset of $\CC$, and $\CC\setminus S$ is an open $0$-neighbourhood. Furthermore, $\ker (I+ Y(\lambda))|_{C^\infty_K}$ is finite-dimensional, and equal to
 		$\ker (I+ \hat{Y}_s(\lambda))$ for each $s\ge s_*$.
 		
 		\item[$(b)$] For all $\lambda\in \CC\setminus S$ and all $s\ge s_*$, the map $I+ \hat{Y}_{s}(\lambda)$ is continuously invertible for all $s\ge s_*$, and $\lambda\mapsto \big(I+ \hat{Y}_{s}(\lambda)\big)^{-1}$ is holomorphic in $\CC\setminus S$. Moreover,
 		$\big(I+ \hat{Y}_{s}(\lambda)\big)^{-1}$ is the restriction of $\big(I+ \hat{Y}_{s_*}(\lambda)\big)^{-1}$ to $H^s_K$ for each $s\ge s_*$, $\lambda\in\CC\setminus S$.
 		
 		\item[$(c)$] $(i)$ Suppose that $f\colon \CC\setminus S\to H^{s_*}_K$ is holomorphic and
 		there is a compact subset ${K_f\subset K}$ so that $\supp \big(\hat{Y}_{s_*}(\lambda)\big)^r f(\lambda)\subset K_f$ for all $r\in\NN_0$ and $\lambda\in \CC\setminus S$. Then $\supp \big(I+\hat{Y}_{s_*}(\lambda)\big)^{-1}f(\lambda)\subset K_f$ for all $\lambda\in \CC\setminus S$.
 		
 		$(ii)$ If $Y(\lambda)C^\infty_K\subset C^\infty_{K'}$ for some compact $K'\subset K$, then
 		$\supp \big(I+ \hat{Y}_{s_*}(\lambda)\big)^{-1}f\subset K'$ for all $\lambda\in \CC\setminus S$ and all $f\in C^\infty_K$, and
 		$\ker \big(I+ \hat{Y}_{s_*}(\lambda)\big)$ is a finite-dimensional subspace of
 		$C^\infty_{K'}$ for all $\lambda\in\CC$.
 		
 		NB. In $(c)$, `support' is to be understood as distributional support.
 	\end{enumerate}
 \end{Theorem}
 \begin{proof}
 	$(a)$ Compactness of $\hat{Y}_s(\lambda)$ follows from the Sobolev embedding theorems.
 	The function $\lambda\mapsto I+ \hat{Y}_{s_*}(\lambda)$ is an analytic function on $\CC$ with values in the Fredholm operators on $H^{s_*}_K$, and which is invertible for $\lambda=0$ because $\hat{Y}_{s_*}(0)=0$. By the analytic Fredholm theorem~\cite[Theorem VI.14]{ReedSimon:vol1}, $I+ \hat{Y}_{s_*}(\lambda)$ is invertible for all $\lambda\in\CC$ with the exception of the (possibly empty) set $S_{s_*}$ of $\lambda\in \CC$ for which $\ker \big(I+ \hat{Y}_{s_*}(\lambda)\big)$ is nontrivial (and necessarily finite-dimensional by~\cite[Theorem VI.15]{ReedSimon:vol1}). Furthermore, $S_{s_*}$ is a
 	discrete subset of $\CC$, whose complement is an open $0$-neighbourhood, and the inverse $\big(I+ \hat{Y}_{s_*}(\lambda)\big)^{-1}$ is meromorphic on $\CC$ and holomorphic on $\CC\setminus S_{s_*}$.
 	
 	If $f\in\ker \big(I+\hat{Y}_{s_*}(\lambda)\big)$ then $f = - \hat{Y}_{s_*}(\lambda)f\in H^{s_*+\delta}_K$ obeys $f\in\ker \big(I+ \hat{Y}_{s_*+\delta}(\lambda)\big)$; iterating, $f\in \bigcap_{s\ge s_*} H^s_K=C^\infty_K$ and $f\in \ker (I+ Y(\lambda))|_{C^\infty_K}$.
 	As the converse inclusion is trivial, we deduce that
 	$\ker \big(I+\hat{Y}_{s_*}(\lambda)\big)= \ker (I+ Y(\lambda))|_{C^\infty_K}$
 	for all $\lambda\in \CC$ and hence that $S_{s_*}=S$ defined in~\eqref{eq:Sdef}. In particular, for each fixed $\lambda\in\CC$, all the kernels $\ker \big(I+\hat{Y}_{s}(\lambda)\big)$ are equal for $s\ge s_*$.
 	
 	$(b)$ The first statement is immediate by the Fredholm theorem and part~(a); the second follows because $I+\hat{Y}_{s}(\lambda)$ coincides with $I+ \hat{Y}_{s_*}(\lambda)$ on $H^s_K$ and both operators are invertible for $\lambda\in\CC\setminus S$.
 	
 	$(c)$\,$(i)$ Let $\chi\in C_0^\infty$ vanish on $K_f$. Regarding each $\big(\hat{Y}_{s_*}(\lambda)\big)^{n}f(\lambda)$ as a distribution with support in $K_f$, we see (e.g., by \cite[Theorem 2.3.3]{Hormander1}) that
 	$\big(\hat{Y}_{s_*}(\lambda)^{n}f(\lambda)\big)(\chi)=0$ for all $n\in \NN_0$ and so the
 	holomorphic function on $\CC\setminus S$ defined by $\lambda\mapsto \big(\big(I+\hat{Y}_{s_*}(\lambda)\big)^{-1}f(\lambda)\big)(\chi)$ vanishes in a neighbourhood of the origin on which the resolvent may be expanded as a convergent geometric series in powers of $\hat{Y}_{s_*}(\lambda)$, recalling that $\hat{Y}_{s_*}(\lambda)\to 0$ as $\lambda\to 0$.
 	By holomorphicity, it follows that
 	$\big(\big(I+\hat{Y}_{s_*}(\lambda)\big)^{-1}f(\lambda)\big)(\chi)=0$ for all $\lambda\in\CC\setminus S$.
 	Allowing $\chi$ to vary, we conclude that $\supp \big(I+\hat{Y}_{s_*}(\lambda)\big)^{-1}f(\lambda)\subset K_f$.
 	
 	$(c)$\,$(ii)$ The assumption implies that of $(c)$\,$(i)$ for all $f\in C^\infty_K$ \big(regarded as constant functions from $\CC\setminus S$ to $H^{s_*}_K$\big) with $K_f=K'$, so
 	$\supp \big(I+\hat{Y}_{s_*}(\lambda)\big)^{-1}f\subset K'$ for all $\lambda\in \CC\setminus S$ and all $f\in C^\infty_K$. Finally,
 	if $g\in \ker \big(I+\hat{Y}_{s_*}(\lambda)\big)=\ker (I+ Y(\lambda))|_{C^\infty_K}$ (already known to be finite-dimensional) then $g\in C^\infty_K$ and $g=- Y(\lambda) g\in C^\infty_{K'}$.
 \end{proof}

 The hypotheses of Theorem~\ref{thm:Sobolevinversion} entail that each $Y(\lambda)$ restricts to a continuous endomorphism of $C^\infty_K$. To see this, note that for any $k\in\NN_0$ and any integer $s>k+n/2$, where $n$ is the maximum dimension of any component of $M$, there are continuous maps
 \begin{equation*}
 	C^s_K \longrightarrow H^s_K \stackrel{Y(\lambda)}{\longrightarrow} H^s_K \longrightarrow C^k_K,
 \end{equation*}
 where the unlabelled arrows are Sobolev embeddings. Thus for all $k\in \NN_0$ there is $s\in \NN_0$ and a constant $C$ such that
 $\|Y(\lambda)f\|_{K,k}\le C \|f\|_{K,s}$ for all $f\in C^\infty_K$. This observation allows the following conclusion to be drawn.
 \begin{Corollary}\label{cor:CKinfty_inversion}
 	In the notation of Theorem~$\ref{thm:Sobolevinversion}$, and for $\lambda\in\CC\setminus S$, $I+ Y(\lambda)$ restricts to a~homeomorphism of $C^\infty_K$, with inverse to be denoted $R(\lambda)$. The map $\lambda\mapsto R(\lambda)$ is
 	holomorphic on $\CC\setminus S$ in the topology of $\LL_b\big(C^\infty_K\big)$.
 	If, additionally, for some holomorphic function $f\colon \CC\setminus S\to C^\infty_K$
 	there is a compact subset $K_f\subset K$ so that $\supp Y(\lambda)^r f(\lambda)\subset K_f$ for all $r\in\NN_0$ and $\lambda\in \CC\setminus S$, then $\supp R(\lambda) f(\lambda)\subset K_f$ for all $\lambda\in \CC\setminus S$. If $Y(\lambda)C^\infty_K\subset C^\infty_{K'}$ for some compact $K'\subset K$, then
 	$\ker(I+ Y(\lambda))|_{C^\infty_K}$ is a finite-dimensional subspace of $C^\infty_{K'}$
 	for all $\lambda\in\CC$.
 \end{Corollary}
 \begin{proof}
 	Suppose $f\in C^\infty_K$ and $\lambda\notin S$. Then, for all $s\ge s_*$,
 	we have	$f\in H^s_K$ and $\big(I+ \hat{Y}_{s}(\lambda)\big)^{-1}f=\big(I+\hat{Y}_{s_*}(\lambda)\big)^{-1}f$, and consequently
 	\begin{equation*}
 		\big(I+\hat{Y}_{s_*}(\lambda)\big)^{-1}f\in \bigcap_{s\ge s_*} H^s_K = C^\infty_K.
 	\end{equation*}
 	Thus $\big(I+\hat{Y}_{s_*}(\lambda)\big)^{-1}$ restricts to linear self-map of $C^\infty_K$ which
 	\big(because $Y(\lambda)$ and $\hat{Y}_{s_*}(\lambda)$ agree on $C^\infty_K$\big) is a linear inverse to the restriction of $I+ Y(\lambda)$. Accordingly, $I+ Y(\lambda)$ restricts to a~continuous bijection of $C^\infty_K$, and since the latter is a Fr\'echet space the inverse mapping theorem implies that the inverse $R(\lambda)$ is continuous.
 	As every $C^\infty_K$ semi-norm is dominated by a~Sobolev norm and
 	vice versa, convergence in the operator norm of every $H^s_K$, for $s\ge s_*$,
 	implies convergence in $\LL_b\big(C^\infty_K\big)$, by Corollary~\ref{cor:HstoCinfty}. It follows that $\lambda\mapsto R(\lambda)$ is holomorphic on~$\CC\setminus S$ in the topology of $\LL_b\big(C^\infty_K\big)$.
 	The remaining statements are immediate from parts~$(a)$ and $(c)$ of Theorem~\ref{thm:Sobolevinversion} and the fact that distributional support of a smooth function is exactly its usual support.
 \end{proof}


 We need the following elementary observation, which will be used for $F=C_K^\infty$, $G=C_0^\infty$.
 \begin{Lemma}\label{lem:inversetransfer}
 	Let $F$ and $G$ be topological vector spaces and suppose that the diagram
 	\begin{equation*} %\label{eq:lemmacd}
 		\begin{tikzcd}
 			F \arrow[r, "S"]\arrow[d,"\iota"] & F \arrow[d,"\iota"]\\
 			G \arrow[r, "T"]\arrow[ur,"\hat{T}"] & G
 		\end{tikzcd}
 	\end{equation*}
 	of continuous linear maps commutes. If $\id_F+S$ is continuously invertible, then
 	$\id_G+T$ is continuously invertible with inverse
 	\begin{equation*}%\label{eq:invlemform}
 		(\id_G+ T)^{-1} = \id_G-T+\iota S(\id_F+S)^{-1}\hat{T}.
 	\end{equation*}
 \end{Lemma}
 \begin{proof}
 	Assuming that $\id_F+S$ is continuously invertible, we compute
 	\begin{align*}
 		(\id_G+T)\big(\id_G-T+\iota S(\id_F+S)^{-1}\hat{T}\big) = \id_G-T^2 + \iota (\id_F+S)S(\id_F+S)^{-1}\hat{T}=\id_G
 	\end{align*}
 	and
 	\begin{align*}
 		\big(\id_G-T+\iota S(\id_F+S)^{-1}\hat{T}\big)(\id_G+T) &= \id_G-T^2 + \iota S(\id_F+S)^{-1}(\id_F+S)\hat{T}= \id_G.
 	\end{align*}
 	using in both cases the identities $\iota S=T\iota$ and hence $\iota S \hat{T} = T\iota\hat{T}=T^2$, and also $\hat{T}T = \hat{T}\iota\hat{T} = S\hat{T}$. Therefore $\id_G+T$ has a linear inverse, given by a manifestly continuous expression.
 \end{proof}

 We come to the proof of the main result.
 \begin{proof}[Proof of Theorem~\ref{thm:GHpert}] The proof involves several steps, and uses some technical lemmas from the Appendix.
 	
 	\emph{$1.$ Preliminary observations.} Define $Y^\pm(\lambda)\in\LL(C_0^\infty,C^\infty)$ by
 	\begin{equation*}
 		Y^\pm(\lambda)=A(\lambda)E^\pm.
 	\end{equation*}
 	Owing to assumptions $(a)$ and $(b)$, we may extend $Y^\pm(\lambda)$ to linear maps $\EE'\to\DD'_K$ with continuous restrictions mapping $H^s_0\to H^{s+\delta}_K$ for all $s\ge s_*$; hence there are also continuous restrictions $Y_s^\pm(\lambda)\colon H^s_K\to H^{s+\delta}_K$ for $s\ge s_*$,
 	and by the Sobolev embedding theorems, compact maps $\hat{Y}^\pm_s(\lambda)\colon H^s_K\to H^{s}_K$ for $s\ge s_*$ which are holomorphic with respect to the operator norm topology, as noted after the statement of Theorem~\ref{thm:GHpert}.
 	
 	In fact, if $K'$ is any compact topologically regular set, the same argument shows that~$Y^\pm(\lambda)$
 define compact maps depending holomorphically on $\lambda$ in the topology of
 	$\LL_b\big(H^s_{K\cup K'}\big)$ for all~${s\!\ge\!s_*}$. By Corollary~\ref{cor:HstoCinfty}, it follows that
 	$Y^\pm(\lambda)$ are holomorphic in $\lambda$ with respect to the $\LL_b\big(C^\infty_{K\cup K'}\big)$ topology.
 	As $C^\infty_K$ and $C^\infty_{K'}$ are continuously embedded topological subspaces of $C^\infty_{K\cup K'}$, and $\Ran Y^\pm(\lambda)\subset C^\infty_K$, we may use Lemma~\ref{lem:Lb_basic}\,$(a)$,\,$(c)$ to deduce that
 	the $Y^\pm(\lambda)$ are holomorphic in $\lambda$ with respect to the
 	topology of $\LL_b\big(C^\infty_{K'},C^\infty_{K}\big)$.
 	By Lemma~\ref{lem:LbLF}, it follows that
 	$A(\lambda)E^\pm$ are also holomorphic with respect to $\LL_b\big(C_0^\infty,C^\infty_K\big)$ and $\LL_b(C_0^\infty)$; using Lemma~\ref{lem:Lb_basic}\,$(a)$, we also have holomorphicity with respect to
 	$\LL_b\big(C^\infty_K,C_0^\infty\big)$.
 	
 	\emph{$2.$ Finite-dimensionality of $\ker (P+ A(\lambda))|_{C^\infty_{\pc/\fc}}$.} Next, observe that $P$ induces bijections between $\ker (P+ A(\lambda))|_{C^\infty_{\pc/\fc}}$ and $\ker \big(I+ Y^{\pm}(\lambda)\big)\big|_{C^\infty_K}$,
 	with inverses given by the restrictions of~$E^{\pm}$. For
 	if $(P+ A(\lambda))\phi=0$ with $\phi\in C^\infty_{\pc/\fc}$ then
 	$P\phi=- A(\lambda)\phi\in C^\infty_K$ and $\phi=E^\pm P\phi$, so $P\phi\in\ker\big(I+ A(\lambda)E^\pm\big)\big|_{C^\infty_K}$; conversely, if $\big(I+ A(\lambda)E^\pm\big)h=0$ with
 	$h\in C^\infty_K$, then $PE^\pm h = h=- A(\lambda)E^\pm h$, so
 	$E^\pm h\in \ker(P+ A(\lambda))|_{C^\infty_{\pc/\fc}}$.
 	Thus
 	\begin{equation*}
 		S^\pm:= \big\{\lambda\in \CC\colon \ker (P+ A(\lambda))|_{C^\infty_{\pc/\fc}}\neq 0\big\}= \big\{\lambda\in \CC\colon \ker \big(I+ Y^{\pm}(\lambda)\big)\big|_{C^\infty_K}\neq 0\big\}.
 	\end{equation*}
 	In combination with Theorem~\ref{thm:Sobolevinversion}\,$(a)$, we also have
 	\begin{equation}\label{eq:GHpert_remark}
 		\dim \ker (P+A(\lambda))|_{C^\infty_{\pc/\fc}} = \dim \ker \big(I+Y_s^{\pm}(\lambda)\big)\big|_{C^\infty_{K}} = \dim \ker \big(I+\hat{Y}_s^{\pm}(\lambda)\big)<\infty
 	\end{equation}
 	\big(the latter kernel taken in $H^s_K$\big) for all $s\ge s_*$. Part~$(A)$ of the Theorem is thus proved.
 	
 	\emph{$3.$ Construction of holomorphic candidate $K$-nonlocal Green operators.} Applying Theorem~\ref{thm:Sobolevinversion} and Corollary~\ref{cor:CKinfty_inversion} to $Y^\pm(\lambda)$, one finds that $S^\pm$ are discrete subsets of $\CC$, whose complements in~$\CC$ are open $0$-neighbourhoods, and the operators $I+ Y^\pm(\lambda)$ are continuously invertible on~$C^\infty_K$ for $\lambda\in\CC\setminus S^\pm$, with inverses
 	that are holomorphic on~$\CC\setminus S^\pm$ in the topology
 	of $\LL_b\big(C^\infty_K\big)$. All these properties hold also for $S=S^+\cup S^-$,
 	so part~(B) is proved.
 	
 	Fixing any $\lambda\in \CC\setminus S$, the diagram of continuous maps
 	\begin{equation*}
 		\begin{tikzcd}
 			C^\infty_K \arrow[r, "T_K"]\arrow[d] & C^\infty_K \arrow[d]\\
 			C_0^\infty \arrow[r, "T_0"]\arrow[ur,"\hat{T}_0"] & C_0^\infty
 		\end{tikzcd}
 	\end{equation*}
 	commutes, where the unlabelled arrows are the canonical inclusions,
 	$T_0$ and $T_K$ are the respective restrictions of $Y^\pm(\lambda)$
 	to $C_0^\infty$, and $C_K^\infty$ and $\hat{T}_0$ exists because
 	$A(\lambda)C^\infty\subset C^\infty_K$.
 	It follows from Lemma~\ref{lem:inversetransfer} that $I+ Y^\pm(\lambda)$ are continuously invertible on $C_0^\infty$. By abuse of notation we write the inverses as $\big(I+ A(\lambda)E^\pm\big)^{-1}$; Lemma~\ref{lem:inversetransfer} now gives the identity
 	\begin{equation}\label{eq:inverseidentity}
 		\big(I+ A(\lambda)E^\pm\big)^{-1} = I - A(\lambda)E^\pm + A(\lambda)E^\pm \big(I+ A(\lambda)E^\pm\big)^{-1} A(\lambda)E^\pm,
 	\end{equation}
 	on $C_0^\infty(K)$, where we suppress notation for inclusions and restrictions. Note that the inverse on the right-hand side is taken in $\LL\big(C^\infty_K\big)$, while the left-hand side is an inverse in $\LL(C^\infty_0)$. Because the former inverse is holomorphic in
 	$\LL_b\big(C^\infty_K\big)$, the Leibniz rule (see Corollary~\ref{cor:Leibniz}) implies that the left-hand side is
 	holomorphic in $\LL_b(C_0^\infty)$; here, we have also used the holomorphicity
 	of~$A(\lambda)E^\pm$ in $\LL_b\big(C_0^\infty,C^\infty_K\big)$, $\LL_b(C_0^\infty)$ and $\LL_b\big(C^\infty_K,C_0^\infty\big)$ established in step 1 of the proof. It follows that the operators
 	\begin{equation*}
 		\tilde{E}^\pm_\lambda= E^\pm \big(I+ A(\lambda)E^\pm\big)^{-1}\in\LL(C_0^\infty, C^\infty),
 	\end{equation*}
 	which are the candidate Green operators for $P+A(\lambda)$,	are holomorphic in $\lambda$ on $\CC\setminus S$ with respect to the topology of $\LL_b(C_0^\infty,C^\infty)$.
 	To prove part (C) it is now enough to check that $\tilde{E}^\pm_\lambda$ are indeed $K$-nonlocal Green operators.
 	
 	\emph{$4.$ Verification that $\tilde{E}^\pm_\lambda$ are $K$-nonlocal Green operators.}
 	Given any $f\in C_0^\infty$,
 	\begin{equation*}
 		g=\big(I+ A(\lambda)E^\pm\big)^{-1} f
 	\end{equation*}
 	is the unique element of $C_0^\infty$ obeying
 	\begin{equation}\label{eq:geqn}
 		g + A(\lambda) E^\pm g= f,
 	\end{equation}
 	whereupon we deduce that $\supp g\subset \supp f \cup K$ and that
 	\begin{equation*}
 		\varphi = E^\pm g = E^\pm \big(I+ A(\lambda)E^\pm\big)^{-1} f = \tilde{E}^\pm_\lambda f
 	\end{equation*}
 	satisfies
 	\begin{equation}\label{eq:Pphi}
 		P\varphi + A(\lambda)\varphi = g+A(\lambda)E^\pm g=f, \qquad \supp \varphi\subset J^\pm(\supp g) \subset J^\pm(K\cup \supp f).
 	\end{equation}
 	In the special case $f=(P+A(\lambda))h$, for $h\in C_0^\infty$,
 	we have
 	\begin{equation*}
 		f = (P+ A(\lambda))E^\pm P h=Ph + A(\lambda) E^\pm Ph
 	\end{equation*}
 	and the unique solution to \eqref{eq:geqn} is clearly $g=Ph$. Thus
 	\begin{equation*}
 		Ph=g=\big(I+ A(\lambda)E^\pm\big)^{-1}f =\big(I+ A(\lambda)E^\pm\big)^{-1}(P+ A(\lambda))h;
 	\end{equation*}
 	consequently $E^\pm \big(I+ A(\lambda)E^\pm\big)^{-1} (P+ A(\lambda))h = h$.
 	In combination with~\eqref{eq:Pphi}, we have shown
 	\begin{equation*}
 		(P+ A(\lambda))\tilde{E}^\pm_\lambda f=f = \tilde{E}^\pm_\lambda (P+ A(\lambda)) f,\qquad\text{and}\qquad
 		\supp\tilde{E}^\pm_\lambda f\subset J^\pm(K\cup \supp f)
 	\end{equation*}
 	for all $f\in C_0^\infty$.
 	Now suppose more specifically that $J^\pm(\supp f)\cap K=\varnothing$ for $f\in C^\infty_0$. Then $A(\lambda)E^\pm f=0$, due to assumption $(d)$, and~\eqref{eq:geqn} is solved by $g=f$, so
 	$\tilde{E}^\pm_\lambda f=E^\pm f$ has support contained in $J^\pm(\supp f)$. Accordingly,
 	\begin{equation*}
 		\supp\tilde{E}^\pm_\lambda f\subset
 		\begin{cases}
 			J^\pm(\supp f), & J^\pm(\supp f)\cap K=\varnothing ,\\
 			J^\pm(K\cup \supp f), & \text{otherwise,}
 		\end{cases}
 	\end{equation*}
 	so $\tilde{E}^\pm_\lambda$ are $K$-nonlocal Green operators for $P+ A(\lambda)$. Part $(C)$ is complete.
 	
 	\emph{$5.$ Continuous extension.} Due to~\eqref{eq:inverseidentity}, one has the formula
 	\begin{equation*}
 		\tilde{E}_\lambda^\pm = E^\pm - E^\pm A(\lambda)E^\pm + E^\pm A(\lambda)E^\pm \big(I+ A(\lambda)E^\pm\big)^{-1}A(\lambda)E^\pm,
 	\end{equation*}
 	in which the three terms on the right-hand side have continuous
 	extensions as maps from $H^s_0$ to $H^{s+\beta}_\loc$, $H^{s+\beta+\delta}_\loc$ and $H^{s+\beta+2\delta}_\loc$ respectively. Because $\delta>0$, we deduce that $\tilde{E}^\pm_\lambda$
 	extends continuously to a map $H^s_0\to H_\loc^{s+\beta}$ as required.
 	By the Leibniz rule (Corollary~\ref{cor:Leibniz}), this extension is holomorphic in $\lambda$ with respect to the topology of $\LL_b\big(H^s_0,H_\loc^{s+\beta}\big)$.
 	This proves part (D).
 	
 	
 	\emph{$6.$ Support non-increasing modifications.} Finally, suppose condition $(d')$ holds, so that one has $\supp A(\lambda)f\subset \supp f$ as well as
 	$\supp A(\lambda)f\subset K$ for all $f\in C^\infty$. Then $(d)$ also holds,
 	either as a consequence of Peetre's theorem~\cite{Peetre:1960} or by the
 	following direct argument:
 	if $f\in C^\infty$ vanishes identically on $K$, then the carrier of $A(\lambda)f\in C^\infty_K$ satisfies
 	both
 	\begin{equation*}
 		\carr A(\lambda)f \subset \inte(K)
 	\end{equation*}
 	and
 	\begin{equation*}
 		\carr A(\lambda)f\subset \supp A(\lambda)f\subset \supp f=\overline{\carr f}\subset \overline{M\setminus K} = M\setminus \inte(K)
 	\end{equation*}
 	and we deduce that $\carr A(\lambda)f$ is empty, i.e., $A(\lambda)f$ vanishes identically. Therefore, all the conclusions reached previously still hold.
 	
 	It further follows that
 	\begin{equation*}
 		\supp Y^\pm(\lambda) f =\supp A(\lambda)E^\pm f\subset K\cap \supp E^\pm f\subset K\cap J^\pm(\supp f)
 		\qquad \text{for all} \ f\in C_0^\infty;
 	\end{equation*}
iterating, we have in particular that
 	\begin{equation*}
 		\supp \big(Y^\pm(\lambda)\big)^r f\subset K^\pm_f:= K\cap J^\pm(\supp f) \qquad \text{for all} \ r\in \NN,\ f\in C_0^\infty \ \text{and} \ \lambda\in\CC\setminus S.
 	\end{equation*}	
 	Therefore $\big(Y^\pm(\lambda)\big)^r Y^\pm(\lambda) f\in C^\infty_{K^\pm_f}\subset C^\infty_K$ for all $r\in\NN_0$, $f\in C_0^\infty$, $\lambda\in \CC\setminus S$.
 	Using Corollary~\ref{cor:CKinfty_inversion}, applied to the function $\lambda\mapsto Y^\pm(\lambda) f\in C^\infty_K$, it follows that both
 	$\supp \big(I+Y^\pm(\lambda)\big)^{-1} Y^\pm(\lambda) f$ and
 	\begin{equation*}
 		\supp Y^\pm(\lambda) \big(I+Y^\pm(\lambda)\big)^{-1} Y^\pm(\lambda) f= \supp \big(I-(I+ Y^\pm(\lambda))^{-1}\big) Y^\pm(\lambda) f
 	\end{equation*}
 	are contained in $K^\pm_f$ for all $\lambda\in \CC\setminus S$. Using the identity~\eqref{eq:inverseidentity},
 	it now follows that
 	\begin{equation*}
 		\supp \big(I+ A(\lambda)E^\pm\big)^{-1} f \subset \supp f\cup \left(K\cap J^\pm(\supp f)\right) \subset
 		J^\pm(\supp f)
 	\end{equation*}
 	and hence $\supp \tilde{E}^\pm_\lambda f \subset J^\pm\big(J^\pm(\supp f)\big)=J^\pm(\supp f)$, which is the support property (G3) for standard Green operators, completing the proof of part $(E)$ and therefore the whole theorem.
 \end{proof}
